% Directly authored alongside the Markdown companion.
% latexmk -pdf -outdir=/tmp/gamma-pi-latex docs/gamma_optimization_pi_brief.tex
\documentclass[10pt]{article}
\usepackage[T1]{fontenc}
\usepackage{amsmath,amssymb}
\usepackage{newtxtext,newtxmath}
\usepackage[letterpaper,margin=0.65in]{geometry}
\usepackage{graphicx,microtype,xcolor,caption,fancyhdr}
\usepackage[colorlinks=true,linkcolor=blue!45!black,urlcolor=blue!55!black]{hyperref}
\definecolor{ink}{HTML}{263442}
\DeclareMathOperator{\sech}{sech}
\DeclareMathOperator{\sign}{sign}
\newcommand{\repo}{https://github.com/samlayton99/precision-mlps/blob/experiment/frozen-gamma-probe}
\newcommand{\evidence}{\repo/results/checkpoint_D_optimizers/expD36_frozen_gamma_probe/full_sweep/refinements/gamma_factorized_kernel}
\graphicspath{{results/checkpoint_D_optimizers/expD36_frozen_gamma_probe/full_sweep/refinements/gamma_factorized_kernel/}}
\hypersetup{pdftitle={Gamma controls readout learning speed},
  pdfsubject={Gamma attenuation, kernel predictions, and measured GD/Adam learning curves}}
\setlength{\parindent}{0pt}
\setlength{\parskip}{5pt plus 1pt}
\setlength{\emergencystretch}{2em}
\setlength{\headheight}{13pt}
\captionsetup{font=small,labelfont=bf,justification=raggedright,singlelinecheck=false}
\pagestyle{fancy}
\fancyhf{}
\fancyhead[L]{\small\color{ink}Gamma controls readout learning speed}
\fancyhead[R]{\small PI research brief}
\fancyfoot[C]{\small\thepage}
\renewcommand{\headrulewidth}{0.3pt}

\begin{document}
\thispagestyle{plain}
{\color{ink}\fontsize{22}{25}\selectfont\bfseries Gamma controls readout learning speed\par}
\vspace{3pt}
\textbf{A small-gamma network can represent a target and still take far longer to learn it.}
Gamma smooths the features, suppressing fine spatial variation. This changes the
readout kernel. Once we calculate that kernel, classical gradient-descent theory
predicts the learning curve and its crossing of an accuracy threshold.

Freeze the centers $c_j$ and common slope $\gamma>0$, and train only the readout from zero:
\begin{equation*}
 f_{\gamma,\theta}(x)=b+\sum_{j=1}^W w_j\tanh(\gamma(x-c_j)).
\end{equation*}
\textbf{Where gamma enters.} Each tanh is a sharp step smoothed over a distance
proportional to $1/\gamma$. We show on page 2 that this smoothing multiplies frequency
$\omega$ by
\begin{equation}\label{eq:filter}
 M_\gamma(\omega)=\frac{z}{\sinh z},\qquad
 z=\frac{\pi|\omega|}{2\gamma},\qquad M_\gamma(0)=1.
\end{equation}
For $|\omega|\gg\gamma$, the multiplier is exponentially small. Applying the
smoothing to both factors of a feature inner product gives the gamma-dependent
kernel exactly.

\textbf{Where the learning-time prediction comes from.}
Let $\Phi_\gamma$ be these features, including the bias, evaluated at $m$ training
inputs, and set $K_\gamma=\Phi_\gamma\Phi_\gamma^T/m$.
For target vector $y\ne0$, half mean squared error, and step
$0<\eta_\gamma\le1/\|K_\gamma\|_2$, classical GD theory [1] gives the relative residual
\begin{equation}\label{eq:training}
 E_\gamma(n)=\frac{\|(I-\eta_\gamma K_\gamma)^ny\|_2}{\|y\|_2}.
\end{equation}
The predicted time to 1\% error is simply the first $n$ for which this curve falls
below $0.01$. This calculation uses the kernel and target, without fitting a
training trajectory. The finite center and sample geometry is retained;
Fourier frequencies need not be kernel eigenvectors.

\begin{figure}[h!]
 \centering
 \includegraphics[width=\linewidth]{pi_brief_three_panel.pdf}
 \caption{\textbf{Same target, different access during optimization.}
 The target is $\sin(2\pi x)+\tfrac12\sin(6\pi x)+\tfrac14\sin(10\pi x)$,
 with 559 fixed-center features and 8193 training inputs on $[-1,1]$.
 \textbf{A:} The analytic gamma filter; vertical guides mark the target frequencies.
 \textbf{B:} Lines are kernel predictions; open circles are actual GD checkpoints,
 subsampled for visibility. Diamonds mark measured first hits of 1\%; dotted
 verticals locate predicted crossings. The predicted and measured crossing counts
 coincide for all four gammas.
 \textbf{C:} Adam curves join saved checkpoints under identical settings.
 Dashed horizontals mark 1\%. Both optimizers start at zero;
 GD uses $\eta_\gamma\simeq0.5/\|K_\gamma\|_2$.}
\end{figure}

For this target, gamma 8 takes \textbf{15,798,313 GD updates} to reach 1\%, versus
\textbf{16,013} at gamma 64. The predicted delay ratio is \textbf{985.70--987.49};
the measured ratio is \textbf{986.59}. The approximation bounds certify the
threshold crossings shown in B. Both models reach 1\%, so this is an optimization
delay at attainable accuracy. Adam also shows a substantial separation with the
same features (C); its curves are empirical, while \eqref{eq:training} applies to GD.

\clearpage
{\color{ink}\large\bfseries Proof of the gamma-to-kernel statement}\par
Let $\rho_\gamma(t)=\gamma\sech^2(\gamma t)/2$. This density has unit mass.
Splitting at a step's center gives
\begin{equation*}
 \int_{\mathbb R}\rho_\gamma(x-t)\sign(t-c)\,dt
 =2\int_{-\infty}^{x-c}\rho_\gamma(u)\,du-1
 =\tanh(\gamma(x-c)).
\end{equation*}
Define the step kernel using the same centers,
\begin{equation*}
 k_{\rm step}(t,s)=1+\sum_j\sign(t-c_j)\sign(s-c_j).
\end{equation*}
Taking feature inner products gives
\begin{equation*}
\begin{split}
 k_\gamma(x,x')&=1+\sum_j\tanh(\gamma(x-c_j))\tanh(\gamma(x'-c_j))\\
 &=\iint\rho_\gamma(x-t)\rho_\gamma(x'-s)k_{\rm step}(t,s)\,dt\,ds.
\end{split}
\end{equation*}
The bias is preserved by unit mass. Bounded features justify exchanging the
finite sum and integrals. Sampling this identity yields
$(K_\gamma)_{i\ell}=k_\gamma(x_i,x_\ell)/m$.

For the Fourier convention $\widehat\rho(\omega)=\int\rho(t)e^{-i\omega t}\,dt$,
set $a=\omega/(2\gamma)$ and $v=e^{2\gamma t}$. The beta integral gives
\begin{equation*}
 \widehat\rho_\gamma(\omega)=\int_0^\infty\frac{v^{-ia}}{(1+v)^2}\,dv
 =\Gamma(1-ia)\Gamma(1+ia)=\frac{\pi a}{\sinh(\pi a)}=M_\gamma(\omega).
\end{equation*}
Here $\Gamma$ is Euler's gamma function. The final equality follows from
$\Gamma(1+ia)=ia\Gamma(ia)$ and the reflection identity
$\Gamma(ia)\Gamma(1-ia)=\pi/\sin(\pi ia)$; the value at zero is one by continuity.
Convolution therefore multiplies the wave $e^{i\omega x}$ by $M_\gamma(\omega)$
in either kernel argument.

For $z>0$, the derivative of $z/\sinh z$ has numerator $\sinh z-z\cosh z<0$:
it starts at zero and its derivative is $-z\sinh z<0$. Consequently, a cap
$\gamma\le\bar\gamma$ bounds the multiplier by $M_{\bar\gamma}(\omega)$.
Also, $\sinh z\sim e^z/2$ gives $M_\gamma(\omega)\sim2ze^{-z}$.
This proves the explicit attenuation and its dependence on gamma, with no
periodic assumption.\hfill$\square$

Writing the residual as $r_n=y-\Phi_\gamma\theta_n$, the standard GD identity
follows from $r_{n+1}=(I-\eta_\gamma K_\gamma)r_n$ and
$r_0=y$. Diagonalization gives a factor $(1-\eta_\gamma\lambda)^n$ in each
eigenmode of eigenvalue $\lambda$. Thus gamma changes the learning rates through
the kernel, and the target's energy in those directions determines the delay.
The cap bounds attenuation; the numerical delay is specific to the target and geometry.

\bigskip
{\small
\textbf{[1]} Y.~Yao, L.~Rosasco, A.~Caponnetto.
\href{https://yao-lab.github.io/publications/YaoCapRos07_EarlyStop.pdf}{On Early
Stopping in Gradient Descent Learning}. \textit{Constructive Approximation}
26:289--315, 2007, \S3.3.\par
\href{\repo/docs/gamma_optimization_paper_note.pdf}{Full study and additional targets}
\enspace$\cdot$\enspace
\href{\repo/results/checkpoint_D_optimizers/expD36_frozen_gamma_probe/full_sweep/REPORT.md}{Experiment protocol}
\enspace$\cdot$\enspace
\href{\repo/docs/gamma_factorized_readout.md}{Finite-kernel construction and certification}
\enspace$\cdot$\enspace
\href{\evidence/pi_brief_figure_data.json}{Figure data and provenance}}
\end{document}
